 % Project description
\def\projectdescription{
The tuning of particle accelerators is a time-consuming and difficult task that continues
to be performed manually in many cases, reducing the scientific output of some of the
most advanced but also costly research facilities in the world. In the field of autonomous
accelerator tuning, researchers aim to develop methods that assist human operators and
thereby reduce tuning times and extend the experimental capabilities of accelerators. One
of the most promising approaches to solve the autonomous accelerator tuning problem is
the use of Bayesian Optimisation (BO), which is already successfully used in operations.~\cite{roussel2024bo} \

However, the most interesting accelerator tuning tasks are characterised by their high
dimensionality. Despite being more capable at dealing with many dimensions than
other algorithms, BO struggles to solve optimisations problems with more than a few
dozen dimensions. To solve this challenge, researchers have proposed to leverage domain
knowledge by replacing the commonly used but uninformative zero mean prior with a
model of the system under optimisation. So far neural networks have been used for this
purpose as they are differentiable and fast to compute. \cite{boltz2025priormeanmodule} While this work has shown the
advantages of using an informed neural network prior, neural networks come with the
disadvantage that they require large amounts of data to be trained, and that outside of the
training data domain, their behaviour is ill-defined. In recent developments, high-speed
differentiable simulators such as Cheetah \cite{kaiser2024bridging} have been proposed as an alternative for use
as a prior in BO, with the key advantage that they do not require any training data, and
are well-behaved throughout. \cite{kaiser2024bridging} \

The goal of this research project is to take the BO with Cheetah prior approach presented
on a very simple example in \cite{kaiser2024bridging}, and implement it on the closer-to-real world as well as
more complex transverse beam parameter tuning task at the ARES accelerator, which has
already been extensively studied with other autonomous tuning methods. The developed
implementation should then be evaluated according to previous evaluations, to establish
how this improved variation of BO compares to other methods and whether extension to
even more complex accelerator tuning tasks is a promising avenue for future research.
}

 % Project tasks
\def\tasks{
\begin{packedenumerate}
	\item Implement BO with Cheetah prior for tuning the transverse beam parameters on a simulation of the ARES Experimental Area.
	\item Evaluate the BO with Cheetah prior implementation in accordance with the evaluation in \cite{kaiser2024rl}.
	\item Evaluate timing and speed of BO with an uninformed prior vs. BO with a Cheetah informed prior.
	\item Study performance under model uncertainties.
	\item \textbf{(Optional depending on accelerator availability)} Evaluate the developed implementation on the real ARES accelerator in accordance with \cite{kaiser2024rl}.
	\item \textbf{(Optional if time allows)} Setup a general implementation of BO with Cheetah.
\end{packedenumerate}
}
